\section{Conclusion}\label{sec:conclusion}

LLM agents for process operations need not pay for a full simulation of every proposal. Learned risk screens that
combine the observable task state and proposed action with, where informative, the agent's internal signals skipped
16--91\% of validator calls in two domains with five models, with certified failure bounds where the data allowed and
far fewer unsafe executions than random skipping in closed loop. Internal signals, and region-based attention grounding
in particular, help when the agent's failures are failures to retrieve the information that decides its answer, and
help little when the agent retrieves the right information and reasons incorrectly. Two rules make such screens safe
and worthwhile in practice: accept proposals unchecked only on first proposals, with thresholds set by a confidence
bound; and deploy screening where validation is costly, since the screen's cost is fixed while the simulator's grows
with fidelity and horizon. Future work should train screens on retries, certify them on larger sets, and test them
with larger models, real-time constraints and plants with richer fault libraries.
